\subsection{Notation and one-atom likelihood geometry}\label{sec:mv-one-atom-geometry}

The deterministic part of both main theorems is the reduction from mixture
likelihood to a Gaussian-bump field. This subsection derives that KKT
criterion, proves monotonicity in the variance defect, and records the exact
single-contact calculation used later in the simplex regime. We first fix the
notation used throughout the paper.

All asymptotic notation is as $n\to\infty$. For a deterministic positive
sequence $r_n$, the relation $U_n=O_{\Pbb}(r_n)$ means that
$|U_n|/r_n$ is tight, while $U_n=o_{\Pbb}(r_n)$ means that
$U_n/r_n\to0$ in probability. We write $r_n\lesssim s_n$ when
$r_n\le C s_n$ for a constant independent of $n$, and use
$\gtrsim,\asymp,\ll$, and $\gg$ analogously. We write
$U_n\xrightarrow{\Pbb}U$ for convergence in
probability and $U_n\Rightarrow U$ for convergence in distribution. A
positive random sequence satisfies $U_n=n^{o_{\Pbb}(1)}$ when, for every
$\eta>0$, $\Pbb\{n^{-\eta}\le U_n\le n^\eta\}\to1$. A
subscript on asymptotic notation records allowed dependence: for example,
$O_{\Pbb,R}(1)$ is tight for each fixed $R$. In iterated limits,
$o_A(1)$ denotes a deterministic quantity that tends to zero as
$A\to\infty$ after the $n\to\infty$ limit. Constants implicit in
$O_m(\cdot)$ may depend on the fixed integer $m$, and other constants may
change from line to line. Point processes carry the vague topology, and a
standard Gumbel variable has distribution function
$s\mapsto\exp\{-e^{-s}\}$.

Let \(d=d_n\ge2\), let \(Z_1,\ldots,Z_n\) be independent
\(\cN(0,I_d)\) vectors, and observe
\[
  X_i=\sqrt{1-\eps_n}\,Z_i,
  \qquad 0\le\eps_n<1.
\]
We fit location mixtures of \(\cN(\mu,I_d)\) densities. The only one-atom
maximum-likelihood candidate is the point mass at the sample mean. Writing
\(\bar Z_n=n^{-1}\sum_iZ_i\) and
\(\widetilde Z_i=Z_i-\bar Z_n\), the directional-derivative criterion says
that this candidate is optimal exactly when
\begin{equation}
\label{eq:one-atom-event}
  E_n^{(1)}
  :=
  \left\{\sup_{t\in\R^d}D_n(t)\le1\right\},
  \qquad
  D_n(t)
  =
  \frac1n\sum_{i=1}^n
  \exp\left\{t\cdot\widetilde Z_i-\frac{\gamma_n|t|^2}{2}\right\},
\end{equation}
where
\begin{equation}
\label{eq:gamma-definition}
  \kappa_n=\frac{\eps_n}{2(1-\eps_n)},
  \qquad
  \gamma_n=\frac1{1-\eps_n}=1+2\kappa_n.
\end{equation}
Indeed, if \(a\) denotes displacement of a fitted component mean from the
sample mean, the usual directional derivative is
\[
  \frac1n\sum_i
  \exp\left\{a\cdot(X_i-\bar X_n)-\frac{|a|^2}{2}\right\}.
\]
The change of variables \(t=\sqrt{1-\eps_n}\,a\) gives
\eqref{eq:one-atom-event}. Notice that \(D_n(0)=1\), so the issue is whether
the field rises above its value at the origin.

Although a multivariate mixture NPMLE need not be unique, this causes no
ambiguity for the event studied here.  The following dimension-free fact
shows that \(E_n^{(1)}\) is almost surely the event that the NPMLE is unique
and one-atomic.

\begin{proposition}
\label{prop:one-atom-unique}
Let \(n\ge2\), and let \(X_1,\ldots,X_n\) be i.i.d.\ from a nondegenerate
spherical Gaussian distribution in \(\R^d\). The probability that
\(\delta_{\bar X_n}\) and a different mixing measure both maximize the
Gaussian-mixture likelihood is zero.
\end{proposition}

\begin{proof}
If \(\delta_{\bar X_n}\) and \(G\ne\delta_{\bar X_n}\) both maximize the
likelihood, strict concavity in the vector of fitted values implies
\(q_G(X_i)=q_{\delta_{\bar X_n}}(X_i)\) for every \(i\).  Since the one-atom
fit is optimal, its physical-coordinate directional field
\[
 \mathcal D_{X,n}(a)
 =\frac1n\sum_{i=1}^n
 \exp\left\{a\cdot(X_i-\bar X_n)-\frac{|a|^2}{2}\right\}
\]
satisfies \(\mathcal D_{X,n}(a)\le1\).
Moreover,
\[
 \int \mathcal D_{X,n}(u-\bar X_n)\,dG(u)
 =\frac1n\sum_i\frac{q_G(X_i)}
 {q_{\delta_{\bar X_n}}(X_i)}=1.
\]
Thus \(\mathcal D_{X,n}(a)=1\) for some \(a\ne0\).

Write \(Y=(X_1-\bar X_n,\ldots,X_n-\bar X_n)=RU\) in polar form in the subspace
\(\{(y_1,\ldots,y_n):\sum_i y_i=0\}\). Conditional on its direction \(U\),
the radius \(R\) has a density. After setting \(t=Ra\), the
directional-derivative field is
\[
 \mathcal D_{R,U}(t)
 =\frac1n\sum_i
 \exp\left\{t\cdot U_i-\frac{|t|^2}{2R^2}\right\}.
\]
For fixed \(U\) and \(t\ne0\), this is strictly increasing in \(R\).
Consequently, for each \(U\), at most one radius can admit both a one-atom
maximizer and a nonzero contact: at any larger radius that same contact
would make the directional derivative exceed one. Thus, conditional on
\(U\), the radius belongs to an at-most-singleton exceptional set, which has
probability zero because \(R\mid U\) has a density. Integration over \(U\)
proves the claim.
\end{proof}

The one-atom event is monotone in the variance defect. This elementary fact
allows all sharp-window limits below to be converted into laws for a
samplewise critical shrinkage.
\begin{lemma}
\label{lem:variance-monotonicity}
Fix $z_1,\ldots,z_n\in\R^d$, and let
$x_i^{(\eps)}=\sqrt{1-\eps}\,z_i$ for $0\le\eps<1$. If the sample based on
$x_1^{(\eps_1)},\ldots,x_n^{(\eps_1)}$ admits a one-atom NPMLE and
$\eps_2\ge\eps_1$, then the sample based on
$x_1^{(\eps_2)},\ldots,x_n^{(\eps_2)}$ also admits a one-atom NPMLE.
\end{lemma}

\begin{proof}
Let $\bar z=n^{-1}\sum_i z_i$ and
\[
 K_n(t)=\frac1n\sum_{i=1}^n
 \exp\left\{t\cdot(z_i-\bar z)-\frac{|t|^2}{2}\right\}.
\]
After the change of variables $t=a\sqrt{1-\eps}$, the one-atom directional
field is
\[
 e^{-\kappa(\eps)|t|^2}K_n(t),
 \qquad
 \kappa(\eps)=\frac{\eps}{2(1-\eps)}.
\]
Because $\kappa$ is increasing, the supremum of this field is nonincreasing
in $\eps$. The directional-derivative criterion proves the claim.
\end{proof}

The next lemma shows that the monotone one-atom set is nonempty and closed,
hence has the form \([\eps_*,1)\). Its endpoint \(\eps_*\) is the samplewise
critical defect used in the introduction.
\begin{lemma}
\label{lem:critical-defect-interval}
For every fixed sample $z_1,\ldots,z_n\in\R^d$, the set of
$\eps\in[0,1)$ for which
$\sqrt{1-\eps}\,z_1,\ldots,\sqrt{1-\eps}\,z_n$ admits a one-atom NPMLE is
a nonempty interval of the form $[\eps_*,1)$.
\end{lemma}

\begin{proof}
Let $y_i=z_i-\bar z$ and $M=\max_i|y_i|$. The empirical random variable
taking the values $t\cdot y_i$ has mean zero and range of length at most
$2M|t|$. Hoeffding's lemma gives
\[
 \log\left(\frac1n\sum_i e^{t\cdot y_i-|t|^2/2}\right)
 \le \frac{M^2-1}{2}|t|^2.
\]
Hence the one-atom field
$e^{-\kappa|t|^2}n^{-1}\sum_i e^{t\cdot y_i-|t|^2/2}$ is at most one
for every $t$ once $\kappa\ge(M^2-1)_+/2$, proving nonemptiness. Lemma
\ref{lem:variance-monotonicity} makes the set an upper interval. Finally, if
$\eps_m$ belongs to the set and $\eps_m\to\eps<1$, then the one-atom
criterion holds at $\eps$ by pointwise convergence of its field for every
$t$. The interval is therefore closed in $[0,1)$ and contains its infimum.
\end{proof}

For the spatial analysis, it is useful to rewrite \(D_n\) as a Gaussian
radial-basis field whose centers and heights are determined by the sample.
This representation will allow the radial and angular contributions to be
studied separately.

\begin{lemma}
\label{lem:bump-representation}
Set
\[
  w=\sqrt{\gamma_n}\,t,
  \qquad
  q_i=\frac{\widetilde Z_i}{\sqrt{\gamma_n}},
  \qquad
  h_i=\frac1n\exp\left\{\frac{|q_i|^2}{2}\right\}.
\]
Define the bump-coordinate field
\[
  B_n(w):=D_n(w/\sqrt{\gamma_n}).
\]
Then, exactly,
\begin{equation}
\label{eq:bump-representation}
  B_n(w)
  =
  \sum_{i=1}^n h_i
  \exp\left\{-\frac{|w-q_i|^2}{2}\right\}.
\end{equation}
In particular, the \(i\)th bump has maximum \(h_i\), attained at \(w=q_i\),
and
\begin{equation}
\label{eq:one-point-threshold}
  h_i>1
  \quad\Longleftrightarrow\quad
  \frac{|\widetilde Z_i|^2}{2}>\gamma_n\log n.
\end{equation}
\end{lemma}

\begin{proof}
Completing the square gives
\[
  t\cdot\widetilde Z_i-\frac{\gamma_n|t|^2}{2}
  =
  w\cdot q_i-\frac{|w|^2}{2}
  =
  \frac{|q_i|^2}{2}-\frac{|w-q_i|^2}{2}.
\]
Summing proves \eqref{eq:bump-representation}; maximizing one summand proves
\eqref{eq:one-point-threshold}.
\end{proof}

The bump representation isolates the basic mechanism. A radial extreme creates a
localized bump whose height is close to one. Other observations contribute a
diffuse background, and one-atomicity fails precisely when the full bump field
exceeds one. At the transition the centers capable of producing a crossing have
radius approximately \(\sqrt{2\log n}\) in \(w\)-coordinates, even when the
original dimension is much larger than \(\log n\).

There is also an exact dual formulation which will be decisive in
Section~\ref{sec:explicit-regime}. Let
\(\Delta_n=\{p\in[0,1]^n:\sum_i p_i=1\}\), let
\(\mathbf u_n=(1/n,\ldots,1/n)\), and write
\[
  K(p\Vert \mathbf u_n)=\sum_{i=1}^n p_i\log(np_i)
\]
for relative entropy, with \(0\log0=0\).

\begin{lemma}
\label{lem:entropy-dual}
For every centered collection \(q_1,\ldots,q_n\in\R^d\),
\begin{equation}
\label{eq:entropy-dual}
  \log\sup_{w\in\R^d}B_n(w)
  =
  \sup_{p\in\Delta_n}
  \left\{
    \frac12\left|\sum_{i=1}^np_iq_i\right|^2
    -K(p\Vert \mathbf u_n)
  \right\}.
\end{equation}
The expression on the right is always nonnegative, since it vanishes at
\(p=\mathbf u_n\).
\end{lemma}

\begin{proof}
The Gibbs variational formula gives, for every \(w\),
\[
  \log\left\{\frac1n\sum_i e^{w\cdot q_i}\right\}
  =
  \sup_{p\in\Delta_n}
  \left\{w\cdot\sum_i p_iq_i-K(p\Vert \mathbf u_n)\right\}.
\]
Subtract \(|w|^2/2\), take the supremum over \(w\), and use
\(\sup_w\{w\cdot v-|w|^2/2\}=|v|^2/2\). Finally,
\(\sum_iq_i=0\), so \(p=\mathbf u_n\) gives zero.
\end{proof}

The point-mass choice \(p=e_i\) in \eqref{eq:entropy-dual} gives the
one-bump obstruction \eqref{eq:one-point-threshold}.  A slightly diffuse
competitor gives an exact and strictly stronger finite-sample threshold.

\begin{proposition}
\label{prop:potts-radial-obstruction}
For \(n\ge3\), define
\begin{equation}
\label{eq:potts-radial-threshold}
 R_n^\star
 =
 \frac{(n-1)^2}{n(n-2)}\log(n-1)
 =
 \log n-\frac1n+O\left(\frac{\log n}{n^2}\right).
\end{equation}
If
\[
 \max_i\frac{|q_i|^2}{2}>R_n^\star,
\]
then \(\sup_wB_n(w)>1\). Consequently, \(E_n^{(1)}\) implies
\(\max_i|q_i|^2/2\le R_n^\star\).
\end{proposition}

\begin{proof}
Fix \(i\), and define \(p\in\Delta_n\) by
\[
 p_i=\frac{n-1}{n},
 \qquad
 p_j=\frac1{n(n-1)}\quad(j\ne i).
\]
Since \(\sum_jq_j=0\),
\[
 \sum_jp_jq_j=\frac{n-2}{n-1}q_i,
 \qquad
 K(p\Vert \mathbf u_n)=\frac{n-2}{n}\log(n-1).
\]
The expression in braces in \eqref{eq:entropy-dual} is therefore positive
exactly when
\[
 \frac{|q_i|^2}{2}>
 \frac{(n-1)^2}{n(n-2)}\log(n-1).
\]
The expansion in \eqref{eq:potts-radial-threshold} follows from Taylor
expansion of \(\log(n-1)\).
\end{proof}

\subsection{Gamma extremes and empirical centering}\label{sec:mv-gamma-normalization}

The normalization in Theorem~\ref{thm:main-support-law} is chosen so that the
largest centered radial energies have an order-one Poisson limit. This
subsection proves that limit uniformly over arbitrary dimension sequences and
quantifies the effect of replacing \(Z_i\) by \(Z_i-\bar Z_n\).

Put
\[
  L_n=\log n,
  \qquad
  k_n=\frac{d_n}{2},
\]
and let \(Y_n\) have the \(\Gamma(k_n,1)\) distribution. Write \(f_n\) and
\(\overline F_n\) for its density and upper tail. Define \(b_n\) and \(a_n\)
by
\begin{equation}
\label{eq:b-a-definition}
  n\overline F_n(b_n)=1,
  \qquad
  a_n=\frac{\overline F_n(b_n)}{f_n(b_n)}.
\end{equation}
Thus \(b_n\) is the exact radial \(1-1/n\) quantile and \(a_n\) is the
reciprocal hazard there. Unlike the fixed-dimensional expansion of \(b_n\),
these definitions remain valid without specifying a relation between \(d_n\)
and \(n\).

We first verify that \(a_n\) is the correct extreme-value scale for every
triangular array of Gamma shapes. The proof is included because uniformity in
\(k_n\) is essential here.

\begin{lemma}
\label{lem:gamma-hazard}
For any sequence \(k_n\ge1\), the quantile \(b_n\) in
\eqref{eq:b-a-definition} satisfies
\[
  \frac{b_n-k_n}{\sqrt{k_n}}\longrightarrow\infty.
\]
Moreover,
\begin{equation}
\label{eq:hazard-asymptotic}
  a_n
  =
  \frac{b_n}{b_n-k_n+1}\{1+o(1)\},
\end{equation}
and, uniformly for \(x\) in compact subsets of \(\R\),
\begin{equation}
\label{eq:gamma-tail-ratio}
  \frac{\overline F_n(b_n+a_nx)}{\overline F_n(b_n)}
  \longrightarrow e^{-x}.
\end{equation}
\end{lemma}

\begin{proof}
We may argue along subsequences. If \(k_n\) remains bounded, then
\(b_n\to\infty\), and the first assertion is immediate. If \(k_n\to\infty\),
the central limit theorem gives
\[
  \Pbb\{Y_n>k_n+M\sqrt{k_n}\}\longrightarrow1-\Phi(M)
\]
for every fixed \(M\), where \(\Phi\) is the standard normal distribution
function. Since \(\overline F_n(b_n)=1/n\to0\), eventually
\(b_n>k_n+M\sqrt{k_n}\). Letting \(M\to\infty\) proves the first claim.

For \(y>k_n-1\), changing variables \(z=y+u\) in the Gamma tail gives the
exact identity
\begin{equation}
\label{eq:tail-density-ratio}
  \frac{\overline F_n(y)}{f_n(y)}
  =
  \int_0^\infty
  \exp\left\{-u+(k_n-1)\log\left(1+\frac uy\right)\right\}\,du.
\end{equation}
Set
\(\delta_y=1-(k_n-1)/y=(y-k_n+1)/y\). The inequality
\(\log(1+v)\le v\) shows that the last integral is at most
\(\delta_y^{-1}\). After the substitution \(v=\delta_yu\), its exponent is
\[
  -v-\frac{k_n-1}{2y^2\delta_y^2}v^2
  +O\left(\frac{k_n}{y^3\delta_y^3}v^3\right)
\]
on every bounded \(v\)-interval. At \(y=b_n\),
\[
  \frac{k_n}{b_n^2\delta_{b_n}^2}
  =
  \frac{k_n}{(b_n-k_n+1)^2}
  \longrightarrow0.
\]
Dominated convergence in \eqref{eq:tail-density-ratio}, using the upper
envelope \(e^{-v}\), therefore yields
\(a_n\delta_{b_n}\to1\), which is \eqref{eq:hazard-asymptotic}.

If \(y=b_n+O(a_n)\), then
\[
  \frac{\delta_y}{\delta_{b_n}}=1+o(1),
  \qquad
  \frac{\overline F_n(y)}{f_n(y)}=a_n\{1+o(1)\},
\]
uniformly over the indicated range. The first relation follows by
differentiating \(\delta_y\) and using
\(k_n/(b_n-k_n+1)^2\to0\); the second follows because the dominated-convergence
argument applied to \eqref{eq:tail-density-ratio} is uniform for
\(y=b_n+O(a_n)\). Since
\((\log\overline F_n)'(y)=-f_n(y)/\overline F_n(y)\), integration from
\(b_n\) to \(b_n+a_nx\) gives \eqref{eq:gamma-tail-ratio}.
\end{proof}

The triangular-array maximum therefore has the usual Gumbel law. No
dimension-dependent centering expansion is needed.

\begin{proposition}
\label{prop:gamma-max}
Let \(Y_{n,1},\ldots,Y_{n,n}\) be independent
\(\Gamma(k_n,1)\) variables. Then, for every \(x\in\R\),
\begin{equation}
\label{eq:gumbel-unscaled}
  \Pbb\left\{
    \max_{1\le i\le n}\frac{Y_{n,i}-b_n}{a_n}\le x
  \right\}
  \longrightarrow
  \exp\{-e^{-x}\}.
\end{equation}
More generally, the number of indices satisfying
\(Y_{n,i}>b_n+a_nx\) converges to a Poisson variable with mean \(e^{-x}\).
\end{proposition}

\begin{proof}
By Lemma~\ref{lem:gamma-hazard},
\(n\overline F_n(b_n+a_nx)\to e^{-x}\). The exceedance count is binomial
with success probability asymptotic to \(e^{-x}/n\), proving the Poisson
limit. Its probability of being zero gives \eqref{eq:gumbel-unscaled}.
\end{proof}

Throughout the radial analysis, put
\begin{equation}
\label{eq:residual-variance}
 \sigma_n^2=1-\frac1n.
\end{equation}
The exact transition coordinate and the associated centered radial excesses
are
\begin{equation}
\label{eq:exact-transition-coordinate}
 s_n^\star
 =\frac{\gamma_nR_n^\star/\sigma_n^2-b_n}{a_n},
 \qquad
 \xi_{n,i}^\star
 =\frac{|Z_i-\bar Z_n|^2/(2\sigma_n^2)-b_n}{a_n}-s_n^\star,
 \qquad
 \delta_n^\star=\frac{\sigma_n^2a_n}{\gamma_n}.
\end{equation}
Thus, for \(q_i=(Z_i-\bar Z_n)/\sqrt{\gamma_n}\),
\begin{equation}
\label{eq:exact-support-radius}
 \frac{|q_i|^2}{2}=R_n^\star+\delta_n^\star\xi_{n,i}^\star.
\end{equation}
These exact quantities will be used in every theorem statement. Some proofs
temporarily replace \(R_n^\star\) by \(\log n\); the comparison is recorded
once in Lemma~\ref{lem:exact-old-edge-coordinate}.

The exact Gamma normalization agrees with the moving-shell normalization in
every fixed dimension. The following bridge is used when the bounded- and
growing-dimensional arguments are combined.
\begin{lemma}
\label{lem:fixed-d-normalization-bridge}
Fix \(d\ge3\), put \(k=d/2\), and define
\[
 b_{n,d}=\log n+(k-1)\log\log n-\log\Gamma(k).
\]
Then \(a_n\to1\), \(b_n-b_{n,d}\to0\), and
\(R_n^\star/\sigma_n^2=\log n+o(1)\). If \(s_n^\star=O(1)\), then
\begin{equation}
\label{eq:fixed-d-coordinate-bridge}
 s_n^\star
 =\frac{\gamma_n\log n-b_n}{a_n}+o(1)
 =\frac{\eps_n}{1-\eps_n}\log n
   -(k-1)\log\log n+\log\Gamma(k)+o(1).
\end{equation}
Moreover, uniformly over \(1\le i\le n\),
\begin{equation}
\label{eq:fixed-d-excess-bridge}
 \frac{|Z_i-\bar Z_n|^2/(2\sigma_n^2)-b_n}{a_n}
 =\frac{|Z_i|^2}{2}-b_{n,d}+o_{\Pbb}(1).
\end{equation}
\end{lemma}

\begin{proof}
The fixed-shape Gamma tail expansion gives
\[
 \overline F_n(y)=\frac{y^{k-1}e^{-y}}{\Gamma(k)}\{1+O(y^{-1})\},
 \qquad y\to\infty.
\]
Substitution in \(n\overline F_n(b_n)=1\) yields
\(b_n=b_{n,d}+o(1)\), and Lemma~\ref{lem:gamma-hazard} gives
\(a_n\to1\). The expansion of \(R_n^\star\) in
\eqref{eq:potts-radial-threshold} gives
\(R_n^\star/\sigma_n^2=\log n+o(1)\). Since \(s_n^\star=O(1)\) forces
\(\gamma_n=O(1)\), the first equality in
\eqref{eq:fixed-d-coordinate-bridge} follows. The second uses
\(\gamma_n-1=\eps_n/(1-\eps_n)\).

Finally, \(\max_i|Z_i|=O_{\Pbb}(\sqrt{\log n})\) and
\(|\bar Z_n|=O_{\Pbb}(n^{-1/2})\). Hence
\[
 \max_i\bigl||Z_i-\bar Z_n|^2-|Z_i|^2\bigr|
 =O_{\Pbb}\left(\sqrt{\frac{\log n}{n}}\right)=o_{\Pbb}(1),
\]
while multiplication by \(\sigma_n^{-2}=1+O(n^{-1})\) changes the maximum
energy by \(O_{\Pbb}(\log n/n)=o_{\Pbb}(1)\). Together with
\(a_n\to1\) and \(b_n-b_{n,d}\to0\), this proves
\eqref{eq:fixed-d-excess-bridge}.
\end{proof}

The Karush--Kuhn--Tucker (KKT) field uses centered observations. The next
condition is sufficient
for empirical centering to be negligible on the Gamma extreme scale:
\begin{equation}
\label{eq:centering-condition}
  \frac{\sqrt{b_nd_n/n}}{a_n}\longrightarrow0.
\end{equation}
It is automatic throughout \(d_n=o(\log n)\).  In high dimension the direct
norm bound behind this condition is wasteful; the leave-one-out estimate
below reaches \(d_n\log n=o(n^2)\).

\begin{lemma}
\label{lem:centered-radial-max}
Assume \eqref{eq:centering-condition}. Then
\begin{equation}
\label{eq:centered-uniform-energy-equivalence}
  \max_i
  \left|
    \frac{|Z_i-\bar Z_n|^2-|Z_i|^2}{2a_n}
  \right|
  \longrightarrow_{\Pbb}0.
\end{equation}
In particular,
\begin{equation}
\label{eq:centered-max-equivalence}
  \max_i
  \frac{|Z_i-\bar Z_n|^2/2-b_n}{a_n}
  -
  \max_i
  \frac{|Z_i|^2/2-b_n}{a_n}
  \longrightarrow_{\Pbb}0.
\end{equation}
Consequently, the centered maximum has the Gumbel limit
\eqref{eq:gumbel-unscaled}.
\end{lemma}

\begin{proof}
Proposition~\ref{prop:gamma-max} and tightness imply
\(\max_i|Z_i|^2/2=b_n+O_{\Pbb}(a_n)\). Lemma~\ref{lem:gamma-hazard} also
implies \(a_n=O(b_n)\), so \(\max_i|Z_i|=O_{\Pbb}(\sqrt{b_n})\). Since
\(|\bar Z_n|=O_{\Pbb}(\sqrt{d_n/n})\),
\[
  \max_i
  \left|
    \frac{|Z_i-\bar Z_n|^2-|Z_i|^2}{2a_n}
  \right|
  \le
  \frac{\max_i|Z_i|\,|\bar Z_n|+|\bar Z_n|^2/2}{a_n}
  =o_{\Pbb}(1)
\]
by \eqref{eq:centering-condition}. This proves
\eqref{eq:centered-uniform-energy-equivalence}; taking maxima gives
\eqref{eq:centered-max-equivalence}.
\end{proof}

The norm bound in Lemma~\ref{lem:centered-radial-max} controls
\(Z_i\cdot\bar Z_n\) by the product of two norms and is unnecessarily
wasteful when \(d_n\) is large. A
leave-one-out decomposition retains the cancellation in this inner product.

\begin{lemma}
\label{lem:centered-radial-max-high-d}
Uniformly over \(d_n\ge1\),
\begin{equation}
\label{eq:centered-energy-leave-one-out}
 \max_{1\le i\le n}
 \left|
  \frac{|Z_i-\bar Z_n|^2-|Z_i|^2}{2}
 \right|
 =
 O_{\Pbb}\left\{
  \frac{d_n}{n}
  +\sqrt{\frac{d_nL_n}{n}}
  +\frac{L_n}{\sqrt n}
 \right\}.
\end{equation}
Consequently, if \(L_n=o(d_n)\) and \(d_nL_n=o(n^2)\), then
\eqref{eq:centered-max-equivalence} holds.
\end{lemma}

\begin{proof}
Write
\[
 Z_i\cdot\bar Z_n
 =
 \frac{|Z_i|^2}{n}
 +Z_i\cdot\bar Z_{-i},
 \qquad
 \bar Z_{-i}=\frac1n\sum_{j\ne i}Z_j.
\]
The two vectors in the second inner product are independent, and
\(\bar Z_{-i}\sim\mathcal N(0,(n-1)n^{-2}I_{d_n})\).  Its moment generating
function therefore gives the subexponential bound
\[
 \max_i|Z_i\cdot\bar Z_{-i}|
 =
 O_{\Pbb}\left\{
  \sqrt{\frac{d_nL_n}{n}}+\frac{L_n}{\sqrt n}
 \right\}.
\]
The standard chi-square tail bound and a union bound also give
\[
 \max_i|Z_i|^2
 =
 O_{\Pbb}\{d_n+\sqrt{d_nL_n}+L_n\},
 \qquad
 |\bar Z_n|^2=O_{\Pbb}\{(d_n+1)/n\}.
\]
Substituting these estimates into
\[
 \frac{|Z_i-\bar Z_n|^2-|Z_i|^2}{2}
 =-Z_i\cdot\bar Z_n+\frac{|\bar Z_n|^2}{2}
\]
proves \eqref{eq:centered-energy-leave-one-out}.

In the stated range,
\(a_n\sim\sqrt{d_n/(4L_n)}\).  Dividing the three terms in
\eqref{eq:centered-energy-leave-one-out} by \(a_n\) gives, respectively,
\[
 O\left(\frac{\sqrt{d_nL_n}}n\right),\qquad
 O\left(\frac{L_n}{\sqrt n}\right),\qquad
 O\left(\frac{L_n^{3/2}}{\sqrt{d_nn}}\right),
\]
all of which tend to zero.  This proves
\eqref{eq:centered-max-equivalence}.
\end{proof}

The two norm-based centering lemmas do not use the exact marginal
distribution of a centered residual. In fact,
\[
 \frac{Z_i-\bar Z_n}{\sigma_n}\sim\mathcal N(0,I_{d_n})
\]
for every \(i\).  The residuals are dependent, but their correlations are
weak enough that the extreme exceedance counts still have Poisson limits.
The next proposition uses this exact marginal law to reach dimensions beyond
the range of Lemma~\ref{lem:centered-radial-max-high-d}.

\begin{proposition}
\label{prop:direct-centered-gamma-max}
Suppose
\[
 L_n=o(d_n).
\]
For every fixed \(x\in\mathbb R\),
\begin{equation}
\label{eq:direct-centered-gumbel}
 \Pbb\left\{
  \max_i
  \frac{|Z_i-\bar Z_n|^2/(2\sigma_n^2)-b_n}{a_n}
  \le x
 \right\}
 \longrightarrow e^{-e^{-x}}.
\end{equation}
More generally, the number of indices for which
\[
 \frac{|Z_i-\bar Z_n|^2}{2\sigma_n^2}>b_n+a_nx
\]
converges to \(\operatorname{Pois}(e^{-x})\).
\end{proposition}

\begin{proof}
Fix \(m\ge1\). We prove, uniformly for \(x\) in compact subsets of
\(\mathbb R\), that
\begin{equation}
\label{eq:centered-joint-tail-factorization}
 \Pbb\{Y_1>y_n,\ldots,Y_m>y_n\}
 =
 \{1+o(1)\}\Pbb\{Y_1>y_n\}^m,
 \qquad y_n=b_n+a_nx,
\end{equation}
where
\[
 V_i=\frac{Z_i-\bar Z_n}{\sigma_n},
 \qquad
 Y_i=\frac{|V_i|^2}{2}.
\]
The proposition will then follow by the method of factorial moments.

We first record the Gamma estimates needed below. Put \(k=k_n=d_n/2\),
\(L=L_n=\log n\), and
\[
 h(\delta)=\delta-\log(1+\delta),\qquad \delta>-1.
\]
If \(Y\sim\Gamma(k,1)\) and \(y=k(1+\delta)\), where
\(\delta\to0\) and \(k\delta^2\to\infty\), then
\begin{equation}
\label{eq:centered-gamma-moderate-tail}
 \Pbb\{Y>y\}
 =
 \frac{1+o(1)}{\delta\sqrt{2\pi k}}
 \exp\{-kh(\delta)\}.
\end{equation}
Writing \(f_k\) for the density of \(Y\), Stirling's formula gives
\[
 f_k\{k(1+\delta)\}
 =
 \frac{1+o(1)}{(1+\delta)\sqrt{2\pi k}}
 \exp\{-kh(\delta)\}.
\]
Moreover, the exact identity
\[
 \frac{\Pbb\{Y>y\}}{f_k(y)}
 =
 \int_0^\infty
 \exp\left\{-v+(k-1)\log\left(1+\frac{v}{y}\right)\right\}\,dv
\]
shows that
\begin{equation}
\label{eq:centered-gamma-mills}
 \frac{\Pbb\{Y>y\}}{f_k(y)}
 =
 \frac{y}{y-k+1}\{1+o(1)\}
 =
 \frac{1+\delta}{\delta}\{1+o(1)\}.
\end{equation}
For completeness, set \(\eta=(y-k+1)/y\) and substitute \(s=\eta v\)
in the integral. Its exponent becomes
\[
 -s-\frac{k-1}{2y^2\eta^2}s^2
 +O\left(\frac{k}{y^3\eta^3}s^3\right)
\]
on bounded \(s\)-intervals, while
\(\log(1+v/y)\le v/y\) supplies the dominating function \(e^{-s}\).
Since \(k/(y^2\eta^2)\sim(k\delta^2)^{-1}\to0\), dominated convergence
proves \eqref{eq:centered-gamma-mills}, and hence
\eqref{eq:centered-gamma-moderate-tail}.

Let
\[
 \delta_n=\frac{b_n-k}{k}.
\]
The central limit theorem and \(\Pbb\{Y>b_n\}=n^{-1}\) imply
\(\delta_n\sqrt{k}\to\infty\). The Chernoff bound
\[
 \Pbb\{Y>k(1+\delta)\}\le e^{-kh(\delta)}
\]
gives \(kh(\delta_n)\le L=o(k)\), and therefore \(\delta_n\to0\).
Applying \eqref{eq:centered-gamma-moderate-tail} at \(b_n\) yields
\[
 L
 =
 kh(\delta_n)+\log(\delta_n\sqrt{2\pi k})+o(1).
\]
Because \(h(\delta)\sim\delta^2/2\), this implies
\begin{equation}
\label{eq:centered-gamma-local-scales}
 \delta_n\sim\sqrt{\frac{2L}{k}},
 \qquad
 b_n-k\sim\sqrt{2kL},
 \qquad
 a_n\sim\frac1{\delta_n}\sim\sqrt{\frac{k}{2L}}.
\end{equation}
If \(y_n=b_n+a_nx=k(1+\delta_{n,x})\), then, uniformly for bounded \(x\),
\[
 \frac{\delta_{n,x}}{\delta_n}=1+o(1).
\]
Using \(h'(\delta)=\delta/(1+\delta)\),
\eqref{eq:centered-gamma-moderate-tail}, and
\(a_n\delta_n/(1+\delta_n)\to1\), we obtain
\[
 k\{h(\delta_{n,x})-h(\delta_n)\}
 =
 x+O\left(\frac{a_n^2}{k}\right)+o(1)
 =
 x+o(1).
\]
Consequently,
\begin{equation}
\label{eq:centered-gamma-local-tail-ratio}
 p_n(x):=\Pbb\{Y_i>y_n\}
 =
 \frac{e^{-x}+o(1)}{n}.
\end{equation}

We next compare the joint law of \(Y_1,\ldots,Y_m\) with \(m\)
independent Gamma variables. For each Euclidean coordinate, the Gaussian
vector \((V_1,\ldots,V_m)\) has covariance
\[
 A=A_{n,m}
 =
 \frac n{n-1}I_m-\frac1{n-1}\mathbf1\mathbf1^\top
 =
 I_m+C,
\]
where
\[
 C_{ii}=0,
 \qquad
 C_{ij}=-\frac1{n-1}\quad(i\ne j),
 \qquad
 \|C\|_{\mathrm{op}}+\|C\|_{\mathrm F}=O_m(n^{-1}).
\]
Set
\[
 t=t_n=1-\frac{k}{y_n},
 \qquad
 \beta=\beta_n=1-t=\frac{k}{y_n}.
\]
By \eqref{eq:centered-gamma-local-scales},
\begin{equation}
\label{eq:centered-tilt-scale}
 t\to0,
 \qquad
 t\sqrt{d_n}=(2+o(1))\sqrt{L_n},
 \qquad
 d_nt^2=(4+o(1))L_n.
\end{equation}

Let \(M_{n,m}^{\mathrm{dep}}(t)\) be the joint moment-generating function
of \((Y_1,\ldots,Y_m)\) at the common argument \(t\). The Gaussian
quadratic-form formula gives
\[
 M_{n,m}^{\mathrm{dep}}(t)
 =
 \det(I_m-tA)^{-d_n/2}.
\]
For \(m\) independent copies,
\[
 M_{n,m}^{\mathrm{ind}}(t)=\beta^{-md_n/2}.
\]
Since \(\operatorname{tr}C=0\),
\begin{align}
\label{eq:centered-tilt-normalizer}
 \log\frac{M_{n,m}^{\mathrm{dep}}(t)}
               {M_{n,m}^{\mathrm{ind}}(t)}
 &=
 -\frac{d_n}{2}
 \log\det\left(I_m-\frac{t}{\beta}C\right) \notag\\
 &=
 O_m\left(\frac{d_nt^2}{n^2}\right)
 =
 O_m\left(\frac{L_n}{n^2}\right)
 =
 o(1).
\end{align}

Under the dependent exponential tilt, each coordinate vector is Gaussian
with covariance
\[
 \Sigma^{\mathrm{dep}}
 =
 (A^{-1}-tI_m)^{-1}
 =
 A(I_m-tA)^{-1}.
\]
Under the independent tilt the corresponding covariance is
\[
 \Sigma^{\mathrm{ind}}=\beta^{-1}I_m.
\]
We write \(\E_t^{\mathrm{dep}}\) and \(\E_t^{\mathrm{ind}}\) for expectation
under these two tilted laws, and use \(\E_t^\star\) when the same formula
applies to either one.
Writing \(f_t(a)=a/(1-ta)\) and expanding \(f_t(I_m+C)\) at \(I_m\)
gives
\[
 \Sigma^{\mathrm{dep}}
 =
 \beta^{-1}I_m+\beta^{-2}C+R_{n,m},
 \qquad
 \|R_{n,m}\|_{\mathrm{op}}
 \le C_m t\|C\|_{\mathrm{op}}^2.
\]
In particular,
\begin{equation}
\label{eq:centered-tilted-covariance}
 (\Sigma^{\mathrm{dep}})_{ii}
 =
 \beta^{-1}+O_m(tn^{-2}),
 \qquad
 (\Sigma^{\mathrm{dep}})_{ij}
 =
 -\frac{\beta^{-2}}{n-1}+O_m(tn^{-2})
 \quad(i\ne j).
\end{equation}
The eigenvalues of both tilted covariance matrices therefore lie in a
fixed compact subset of \((0,\infty)\).

We now prove the local limit estimate needed for the tilted overshoots.
Let \(f_{n,m}^{\mathrm{dep}}\) and \(f_{n,m}^{\mathrm{ind}}\) denote the
respective densities of
\[
 W_n
 =
 \frac1{\sqrt{d_n}}
 (Y_1-y_n,\ldots,Y_m-y_n)
\]
under the dependent and independent tilted laws. We claim that
\begin{equation}
\label{eq:centered-tilted-local-limit}
 \sup_{w\in\mathbb R^m}
 \left|f_{n,m}^{\star}(w)-\varphi_m(w)\right|
 \longrightarrow0,
 \qquad
 \varphi_m(w)=\pi^{-m/2}e^{-|w|^2},
\end{equation}
for \(\star\in\{\mathrm{dep},\mathrm{ind}\}\).

To prove the claim, let \(\Sigma^\star\) be the appropriate tilted
covariance, let \(D_z=\operatorname{diag}(z_1,\ldots,z_m)\), and define
\[
 B_z^\star=(\Sigma^\star)^{1/2}D_z(\Sigma^\star)^{1/2}.
\]
The characteristic function of \(W_n\) is
\begin{equation}
\label{eq:centered-tilted-characteristic-function}
 \chi_{n,m}^{\star}(z)
 =
 \exp\left\{-\frac{iy_n}{\sqrt{d_n}}\sum_{j=1}^mz_j\right\}
 \det\left(
 I_m-\frac{i}{\sqrt{d_n}}\Sigma^\star D_z
 \right)^{-d_n/2}.
\end{equation}
For fixed \(z\), expansion of the logarithm gives
\begin{align*}
 \log\chi_{n,m}^{\star}(z)
 &=
 \frac{i\sqrt{d_n}}2
 \sum_{j=1}^m
 \left\{(\Sigma^\star)_{jj}-\beta^{-1}\right\}z_j\\
 &\quad
 -\frac14
 \operatorname{tr}(\Sigma^\star D_z\Sigma^\star D_z)
 +O_m\left(\frac{|z|^3}{\sqrt{d_n}}\right).
\end{align*}
The linear term vanishes identically in the independent case and is
\(O_m(t\sqrt{d_n}|z|/n^2)=o(1)\) in the dependent case by
\eqref{eq:centered-tilt-scale} and
\eqref{eq:centered-tilted-covariance}. Since
\(\Sigma^\star\to I_m\), it follows that
\[
 \chi_{n,m}^{\star}(z)\longrightarrow e^{-|z|^2/4}
\]
for every fixed \(z\).

We next give an integrable domination uniform in the two tilted laws.
Since \(B_z^\star\) is real symmetric,
\begin{equation}
\label{eq:centered-characteristic-modulus}
 |\chi_{n,m}^{\star}(z)|
 =
 \det\left(
 I_m+\frac{(B_z^\star)^2}{d_n}
 \right)^{-d_n/4}.
\end{equation}
If \(\mu_1,\ldots,\mu_m\) are the eigenvalues of \(B_z^\star\), the
uniform lower eigenvalue bound for \(\Sigma^\star\) gives
\[
 \sum_{j=1}^m\mu_j^2
 =
 \|B_z^\star\|_{\mathrm F}^2
 \ge c_m|z|^2.
\]
Hence
\begin{equation}
\label{eq:centered-characteristic-global-bound}
 |\chi_{n,m}^{\star}(z)|
 \le
 \left(1+\frac{c_m|z|^2}{d_n}\right)^{-d_n/4}.
\end{equation}
Choose \(c_0>0\) sufficiently small. If
\(|z|\le c_0\sqrt{d_n}\), then all
\(\mu_j^2/d_n\le1\); using \(\log(1+v)\ge v/2\) in
\eqref{eq:centered-characteristic-modulus} yields
\begin{equation}
\label{eq:centered-characteristic-gaussian-bound}
 |\chi_{n,m}^{\star}(z)|\le e^{-c_m'|z|^2}.
\end{equation}
On the complementary region, polar coordinates and
\eqref{eq:centered-characteristic-global-bound} give
\begin{align*}
 &\int_{|z|>c_0\sqrt{d_n}}
 \left(1+\frac{c_m|z|^2}{d_n}\right)^{-d_n/4}\,dz\\
 &\qquad\le
 C_m d_n^{m/2}(1+c_m'')^{-d_n/8}
 \mathrm B\left(\frac m2,\frac{d_n}{8}-\frac m2\right)
 =o(1).
\end{align*}
Here \(\mathrm B\) denotes the Euler beta function.
Thus the characteristic functions are uniformly integrable. Pointwise
convergence, \eqref{eq:centered-characteristic-gaussian-bound}, and the
last tail estimate imply
\[
 \int_{\mathbb R^m}
 \left|\chi_{n,m}^{\star}(z)-e^{-|z|^2/4}\right|\,dz
 \longrightarrow0.
\]
Fourier inversion proves \eqref{eq:centered-tilted-local-limit}. Notice
that this argument treats the dependent and independent tilted densities
separately; both converge uniformly to the same density \(\varphi_m\), and
in particular
\[
 f_{n,m}^{\mathrm{dep}}(0),
 \ f_{n,m}^{\mathrm{ind}}(0)
 \longrightarrow
 \varphi_m(0)=\pi^{-m/2}.
\]

It remains to compare the tilted overshoot integrals. For
\(\star\in\{\mathrm{dep},\mathrm{ind}\}\), put
\[
 J_{n,m}^{\star}
 =
 \E_t^\star\left[
  \exp\left\{-t\sum_{j=1}^m(Y_j-y_n)\right\}
  \mathbf1_{\{Y_j>y_n,\ 1\le j\le m\}}
 \right].
\]
Writing \(q_n=t\sqrt{d_n}\), the density of \(W_n\) gives
\[
 q_n^mJ_{n,m}^{\star}
 =
 \int_{\mathbb R_+^m}
 e^{-\sum_jz_j}
 f_{n,m}^{\star}(z/q_n)\,dz.
\]
By \eqref{eq:centered-tilt-scale}, \(q_n\to\infty\). The uniform local
limit theorem supplies a common bound on the densities, so dominated
convergence yields
\[
 q_n^mJ_{n,m}^{\star}\longrightarrow\varphi_m(0).
\]
Therefore
\[
 \frac{J_{n,m}^{\mathrm{dep}}}{J_{n,m}^{\mathrm{ind}}}\longrightarrow1.
\]
The exponential-tilting identities are
\[
 \Pbb\{Y_1>y_n,\ldots,Y_m>y_n\}
 =
 M_{n,m}^{\mathrm{dep}}(t)e^{-tmy_n}J_{n,m}^{\mathrm{dep}}
\]
and
\[
 p_n(x)^m
 =
 M_{n,m}^{\mathrm{ind}}(t)e^{-tmy_n}J_{n,m}^{\mathrm{ind}}.
\]
Together with \eqref{eq:centered-tilt-normalizer}, they prove
\eqref{eq:centered-joint-tail-factorization}.

Finally, let
\[
 N_n(x)=\sum_{i=1}^n\mathbf1_{\{Y_i>b_n+a_nx\}}.
\]
For an integer-valued variable \(N\), write
\(\{N\}_m=N(N-1)\cdots(N-m+1)\), and similarly
\((n)_m=n(n-1)\cdots(n-m+1)\).
For every fixed \(m\), exchangeability,
\eqref{eq:centered-joint-tail-factorization}, and
\eqref{eq:centered-gamma-local-tail-ratio} give
\[
 \E\{N_n(x)\}_m
 =
 (n)_m\Pbb\{Y_1>y_n,\ldots,Y_m>y_n\}
 \longrightarrow e^{-mx}
\]
for every fixed \(m\). These are the factorial moments of
\(\operatorname{Pois}(e^{-x})\), so
\(N_n(x)\Rightarrow\operatorname{Pois}(e^{-x})\). Taking the probability
that \(N_n(x)=0\) proves \eqref{eq:direct-centered-gumbel}.
\end{proof}
